\begin{proposition}[The necessary normalization defect in the three named test families]\hypertarget{tag-01FQ}{}\label{cdp:st:C4g}
\StatementRecord{cdp:st:C4g}{proved}{01FQ}{\SeriesRef{cdp:st:E13}, \SeriesRef{cdp:st:E14}, \SeriesRef{cdp:st:A.3.4}, \SeriesRef{cdp:st:A.4.1}, \SeriesRef{cdp:st:C.1.3}, \SeriesRef{cdp:st:C5b}}

For every admissible member of the original $(3,4)$ family,
every member of the two-parameter family in \ref{cdp:st:E13}, and every
member of the Liu--Rollenske family in \ref{cdp:st:E14}, there is a reduced
full hexagonal fibre $W$ with normalization
$\eta:\mathrm dP_6\to W$ such that
\[
 \eta^*L=\OO_{\mathrm dP_6}\quad(L\in\Pic(Y)),\qquad
 H^2\bigl(\mathrm dP_6,K_{\mathrm dP_6}\otimes
                    \eta^*(K_Y^{-1}\otimes L)\bigr)=\C.
 \tag{C.TEST.1}\hypertarget{tag-012O}{}\label{cdp:eq:C.TEST.1}
\]
In these three families the singular-fibre Euler contributions are
\[
\begin{array}{c|c|c}
\text{family}&\text{reduced singular supports}&
                 \text{their Euler numbers}\\ \hline
(3,4)&W&2\\
\text{two-parameter \ref{cdp:st:E13}}&W,\ W_1,\ W_2&2,0,0\\
\text{Liu--Rollenske \ref{cdp:st:E14}}&W,\ C_{\rm cusp}&2,0 .
\end{array}
\tag{C.TEST.2}\hypertarget{tag-012P}{}\label{cdp:eq:C.TEST.2}
\]
The other critical full fibres are smooth multiple fibres,
with multiplicities $(3,4)$, $4$, and $4$, respectively.
Every one of those full multiple fibres passes the generic
ambient vanishing test \ref{cdp:st:A.3.4}, including all its nilpotent
layers. Consequently positive $c_3$ in these examples does
not require a failure of that test on a nonreduced fibre.
The reduced full fibre $W$ already supplies the fixed
normalization obstruction, and $H^2(Y,T_Y\otimes L)\ne0$
for every ambient line bundle in each family.
\end{proposition}


\begin{proof}

\begin{enumerate}[label={$\langle1\rangle$\arabic*.},leftmargin=*]
\item The original local filling is the hexagonal toric quotient
in \ref{cdp:st:A.4.1}. The two-parameter construction in \ref{cdp:st:E13} places this
same filling at the simple nodal value where its specified
linearization vanishes; its other two simple nodal values
have rank-one fillings. The coordinate and line-bundle
comparisons of \ref{cdp:st:E14} place the hexagonal filling at its
separate simple node and the cuspidal quotient at the
remaining critical value. Those proofs construct the actual
local spaces and their identifications over the punctured
base discs. Thus the normalization of $W$ in each case is
the stated smooth rational surface.

\item For every member in question, $H^2(Y,\Z)$ is a finite
group. For the original affine family this follows from E19:
the universal coefficient sequence gives
\[
 H^2(Y,\Z)\simeq\operatorname{Ext}^1(H_1(Y,\Z),\Z)
 \simeq\Z/d\Z,\qquad d=|p|,
\]
because $H_2(Y,\Z)$ is torsion. In the two other families
the integral homology is that of the six-sphere. In
particular every ambient first Chern class is torsion.
The group $H^2(\mathrm dP_6,\Z)$ is free, so its pullback
to $\mathrm dP_6$ is zero. The exponential sequence there,
together with $H^1(\mathrm dP_6,\OO)=0$, shows that the
pulled-back line bundle is trivial. This applies to both
$L$ and $K_Y^{-1}\otimes L$, and Serre duality on
$\mathrm dP_6$ proves \eqref{cdp:eq:C.TEST.1}.

\item All these threefolds have $b_1=b_2=b_3=0$: use \CompendiumRef{II}{cdp:st:E19} for
every original affine member, and the integral homology
of the six-sphere for the other two families. Any full
fibre is a nonzero effective Cartier divisor. Statement
\ref{cdp:st:C.1.3} therefore gives
\[
 h^2(Y,T_Y\otimes L)\geq 1
 \qquad\text{for every }L\in\Pic(Y).
\]
This argument applies equally when $d>1$ and the integral
cohomology has nonzero torsion. It requires no surjectivity
of the global normalization cohomology map. Indeed \ref{cdp:st:C5b}
computes the actual exceptions to that surjectivity in
the original family.

\item The toric strata of $W$ consist of one $(\C^*)^2$ stratum,
three $\C^*$ strata and two triple points. Their compactly
supported Euler numbers give $e(W)=2$.
Each rank-one normalization is a $\PP^1$-bundle over an
elliptic curve, with its zero and infinity sections glued.
Both the normalization and these conductor curves have
Euler number zero, so normalization additivity gives
$e(W_i)=0$. At the cusp in \ref{cdp:st:E14}, normalization of the
cuspidal rational curve makes its degree-zero line bundle
trivial. The lifted contracting action on this normalized
line complement is a product of a fixed automorphism of
$\PP^1$ and multiplication by a constant of modulus
strictly less than one on $\C^*$. Its quotient is a
$\PP^1$-bundle over the elliptic curve given by that
multiplicative period. The cusp normalization is
bijective; a proper continuous bijection between these
compact Hausdorff spaces is a homeomorphism. Hence
$e(C_{\rm cusp})=0$. This proves \eqref{cdp:eq:C.TEST.2}.

\item Finally the explicit good-reduction covers in \ref{cdp:st:E13}--\ref{cdp:st:E14}
and the original root charts have the stated orders.
Their reduced central fibres are smooth finite unramified
torus quotients, and their Euler numbers are zero.
\ref{cdp:st:A.3.4} applies to their actual full fibres, proving the
asserted generic ambient vanishing. Euler localization
therefore gives $c_3(Y)=2$ in all three rows, while the
nonreduced fibres satisfy the complete multiple-fibre
test. This is the required test of the proposed additional
nonreduced-fibre necessity.

\end{enumerate}
\end{proof}